Theorem 8. Let $K$ be a complex, such that $M=|K|$ is a 2 -manifold. Then $K$ is a combinatorial 2-manifold. That is, every subcomplex St $v$ is a combinatorial 2-cell.

Proof as follows.

Lemma 1. Every vertex of $K$ lies in an edge of $K$.\\
Proof. Because $\mathbf{R}^{2}$ has no isolated points. \(\square\)

Lemma 2. Every edge $v_{0} v_{1}$ of $K$ lies in at least one 2 -simplex of $K$.

Proof. Let $v$ be an interior point of $v_{0} v_{1}$, and suppose that $v_{0} v_{1}$ lies in no 2-simplex of $K$. Then $v$ has a neighborhood $U$ in $|K|$, lying in $\operatorname{Int} v_{0} v_{1}$ and homeomorphic to the plane. This is impossible, because no point separates the plane. \(\square\)

Lemma 3. Every edge of $K$ lies in at least two 2 -simplexes of $K$.\\
Proof. Suppose that $e$ lies in only one $\sigma^{2} \in K$; and let $v$ and $U$ be as in the proof of Lemma 2. Then $\sigma^{2}$ contains a semicircle which lies in $U$ and separates $U$, and this is impossible, because no arc separates $\mathbf{R}^{2}$. \(\square\)

Lemma 4. Every edge of $K$ lies in at most two 2 -simplexes of $K$.\\
Proof. Suppose that $e$ is the intersection of three distinct simplexes, $e=\sigma_{1}^{2} \cap \sigma_{2}^{2} \cap \sigma_{3}^{2}$, and let $v$ be an interior point of $e$. Let $U$ be a neighborhood of $v$ in $K$, homeomorphic to $\mathbf{R}^{2}$, and let $\varepsilon>0$ be sufficiently small so that $N(v, \varepsilon) \cap|K| \subset U$. Let $C_{1}$ and $C_{2}$ be semicircles in $\sigma_{1}^{2}$ and $\sigma_{2}^{2}$, with $v$ as center, and the same radius, sufficiently small so as to lie in $N(v, \varepsilon)$. Then $J=C_{1} \cup C_{2}$ is a 1 -sphere in $U$, not separating $U$; and this contradicts the Jordan curve theorem. \(\square\)

We recall, from Section 0, that the link $L(v)$ is the set of all simplexes of St $v$ that do not contain $v$.

Lemma 5. Each set $|L(v)|$ is connected.\\
Proof. If not, $v$ separates $\mid$ St $v \mid$; and this is impossible, because no point separates any open set in $\mathbf{R}^{2}$. \(\square\)

From Lemma 1 it follows that no set $|L(v)|$ is empty. By Lemmas 2, 3, and 4 , every component of $|L(v)|$ is a polygon; and by Lemma 5 it then follows that $|L(v)|$ is a polygon. Now take any $\sigma^{2}$, and subdivide the edges so as to get a 1 -dimensional complex with the same number of edges as $L(v)$. Form a subdivision of $\sigma^{2}$, using the vertices of the subdivision of Fr $\sigma^{2}$, and using one new vertex which lies in no edge of $\sigma^{2}$. There is now a simplicial homeomorphism between the complexes shown in Figures 4.10(a) and 4.10(b).

\begin{center}
\fbox{\parbox{0.82\linewidth}{\centering \emph{Figure 4.10 omitted in this rendered excerpt; see the raw book source above.}}}
\end{center}


